% =====================================================================
%  Appendix: robustness of the wealth label to its rate parameters.
%  Standalone -- \input after \appendix. Self-contained: refers only to
%  Section \ref{sec:wealth} and its equations. booktabs + natbib are
%  loaded by main.tex. Numbers come from the 2011--2023 U.S.-metros
%  panel (54,519--55,061 tracts, all CBSAs > 500k population).
% =====================================================================

\section{Robustness of the Wealth Label to the Discount and Yield Rates}
\label{app:label-robustness}

The occupant-wealth index $W_{i,t}$ of Section~\ref{sec:wealth} depends on two rates: the
real discount rate $r$ that capitalizes the human-capital annuity [Eq.~\eqref{eq:hc}] and the
gross asset yield $r_k$ that capitalizes the financial-asset perpetuity [Eq.~\eqref{eq:cap}].
The main text fixes $r=0.05$ and $r_k=0.045$. This appendix shows that the label the model
trains on---the within-city relative score $Z^{W}_{i,t}$ of Eq.~\eqref{eq:zw}---barely moves
across plausible values of either rate, so the two rates are not free parameters to be tuned.

\paragraph{The label is a near-invariant of the discount rate.}
Because $Z^{W}_{i,t}$ standardizes $\ln W_{i,t}$ within city and year, any rate change that
rescales the index roughly monotonically washes out of the $z$-score. We recompute the label
for each $r$ in $\{0.02, 0.03, 0.05, 0.07\}$ and correlate the resulting label sets across the
full panel. Table~\ref{tab:rate-robust} reports the pairwise Spearman rank correlations.
Across all six rate pairs the rankings never fall below $\rho = 0.988$, with a median of
$\rho = 0.997$; the $r=0.05$ label is at least $\rho = 0.995$ correlated with the label at
every other rate. The pattern holds in the first, middle, and last panel years.

\begin{table}[ht]
\centering
\caption{Pairwise Spearman rank correlation of the relative wealth label $Z^{W}_{i,t}$ across
the discount-rate grid $r\in\{0.02,0.03,0.05,0.07\}$, by year. Pooled over all tracts in all
CBSAs $>500$k population ($n\approx 5.5\times10^{4}$). Each cell summarizes the six pairwise
correlations among the four rate-specific labels.}
\label{tab:rate-robust}
\begin{tabular}{lcccc}
\toprule
Year & Min $\rho$ & Median $\rho$ & Max $\rho$ & vs.\ $r=0.05$, min $\rho$ \\
\midrule
2014 & 0.991 & 0.997 & 0.999 & 0.996 \\
2019 & 0.990 & 0.997 & 0.999 & 0.996 \\
2023 & 0.988 & 0.997 & 0.999 & 0.995 \\
\bottomrule
\end{tabular}
\end{table}

\paragraph{The asset yield does not change the ordering.}
The financial-asset term \eqref{eq:cap} enters \eqref{eq:w2} additively and is small relative
to the human-capital and equity terms, so $r_k$ rescales a minor component without reordering
tracts. Varying $r_k$ over $\{0.040, 0.045, 0.050\}$ changes the Spearman correlation between
$W_{i,t}$ and per-capita income by at most $0.003$ at any discount rate, and by $0.002$ at the
headline $r=0.05$ (Table~\ref{tab:rk-robust}).

\begin{table}[ht]
\centering
\caption{Sensitivity of the wealth index to the gross asset yield $r_k$. Entries are the
spread (max $-$ min) in the Spearman correlation between $W_{i,t}$ and per-capita income as
$r_k$ ranges over $\{0.040,0.045,0.050\}$, by discount rate, 2023.}
\label{tab:rk-robust}
\begin{tabular}{lcccc}
\toprule
Discount rate $r$ & 0.02 & 0.03 & 0.05 & 0.07 \\
\midrule
$\rho(W,\text{income})$ spread over $r_k$ & 0.003 & 0.003 & 0.002 & 0.001 \\
\bottomrule
\end{tabular}
\end{table}

\paragraph{The structural-change classification is robust.}
The rate could in principle matter for the longitudinal test of Eq.~\eqref{eq:sbreak-w},
which flags tracts whose relative wealth moved significantly between 2014 and 2024. It does
not. Running the test separately for each $r$, the indicator $I^{\,\mathrm{valid}}_{i}$ agrees
across all four rates for $95.1\%$ of tracts. The share of tracts flagged ranges from $11.1\%$
at $r=0.02$ to $14.1\%$ at $r=0.07$---a higher rate front-loads the human-capital annuity and
so registers earnings changes a little more readily---with $r=0.05$ at $12.8\%$. No conclusion
of the event study turns on this margin.

\paragraph{Implication.}
The candidate labels are $\rho \geq 0.988$ identical, a gap smaller than the ACS margin of
error embedded in the labels and well below the rank fidelity the vision model can resolve. A
search over $(r,r_k)$ would select among indistinguishable targets. We therefore adopt the
single configuration $(r,r_k)=(0.05,0.045)$ and report the grid here as a robustness check.
